
\onecolumn

\noindent
\textbf{Reviewer DPV5}

\color{blue}
The paper makes many claims about the advantages, specifically, differences from top-down approaches. However, the advantages were not demonstrated empirically in a systematic and scientifically reliable manner.

Experimental results were basically about network properties, which assume the obtained ones are correct. The technical approach and the obtained knowledge were not evaluated or validated.

In addition, the practical impact of the analysis is not so explicit. If the objective of the paper is about network properties and analysis of supply chains, the background, motivation, and definition of research questions should have been given as the central topic of the paper, not in an ad-hoc and incremental way in the experiments description.
\color{black}

\textbf{Response.}
These are good points, and we will clarify all of them in the final version.
We would like to point out that the results are not just about network analysis.
One of our significant results is on Disruption Analysis (section 4.4) using the generated hypergraphs.
This is one of the motivating questions for our work, and shows interesting effects of trade disruptions for different types of commodities.
We will discuss this more in the Introduction.

\noindent
\textbf{Reviewer UUGx}

\color{blue}

    My main concern is an absence of proper correctness evaluation of the extracted hyperedges by using some sort of expert verification or a ground-truth knowledge base. While the pipeline includes rule-based filtering and external tool lookups, it still relies very heavily on LLM-based interpretation. As a result, it is difficult to assess the reliability of the discovered transformations. Even a small-scale expert evaluation or comparison against a subset of known transformations would significantly strengthen the authors' contribution. Furthermore, the topological property and disruption analyses are interesting and valuable, but their interpretability depends on the correctness of the underlying hyperedges. Without some measure of extraction reliability, it is difficult to draw firm conclusions from these analyses.
    The hypergraph representation and the architectural structure are well presented, motivated and executed but not fundamentally novel. The contribution lies mostly in system integration. This is not necessarily a weakness, but it should be acknowledged properly and the authors should properly contextualise their work wrt the literature. Some examples that may be relevant for contextualisation (not exhaustive nor complete list):\\
        Yang, Z., Dong, L., Du, X., Cheng, H., Cambria, E., Liu, X., ... & Wei, F. (2024, March). Language models as inductive reasoners. In Proceedings of the 18th Conference of the European Chapter of the Association for Computational Linguistics (Volume 1: Long Papers) (pp. 209-225).\\
        Shang, W., \& Huang, X. (2025). A survey of large language models on generative graph analytics: Query, learning, and applications. IEEE Transactions on Knowledge and Data Engineering.\\
        Arsenyan, V., Bughdaryan, S., Shaya, F., Small, K. W., & Shahnazaryan, D. (2024, August). Large language models for biomedical knowledge graph construction: information extraction from EMR notes. In Proceedings of the 23rd Workshop on Biomedical Natural Language Processing (pp. 295-317).\\
        Chu, Z., Wang, Y., Cui, Q., Li, L., Chen, W., Qin, Z., & Ren, K. (2024). Llm-guided multi-view hypergraph learning for human-centric explainable recommendation. arXiv preprint arXiv:2401.08217.\\
    Statements regarding domain generality may be better framed as architectural rather than empirical generality, given that only one domain is experimentally evaluated. Otherwise, the authors should justify why the domain used in the paper allows for broad extrapolation.

\color{black}

\textbf{Response}
What is most novel in our work is the mechanism we built where every predicted hyperedge must be supported by explicit references and then scored for evidence alignment. This forces the LLM to act as a kind of fact-checking agent rather than a free-form generator.

The key point is that this produces structured, reference-grounded factual assertions that can actually be evaluated by humans. Unlike top-down supply-chain representations, where the intermediate steps are opaque and unverifiable, our approach makes correctness testable: a human can read the predicted transformation, look at the cited materials, and judge whether the LLM's score is justified.

Thus, the contribution is not ``yet another supply-chain graph,'' but a generalizable, reference-aligned method for validating structured claims generated by LLMs. This makes the work more broadly relevant and scientifically defensible, and it aligns well with what the reviewers found interesting: the hybrid symbolic/LLM fact-checking approach for reliable knowledge generation.

The sensitivity analysis (section 4.3) also helps in understanding robustness of the constructed hypergraphs.
We will clarify this more in the final version.

\noindent
\textbf{Reviewer S2X2}
\color{blue}
The main problem of the paper is that it does not present any contribution to the AAMAS area.

The authors claim that they have used an agentic approach to create a knowledge hypergraph for materials but the architecture or any other detail about the implementation is provided. No detail about how a LLM was used: did you used prompt engineering? did you used a fine-tuned model? did you used any technique of (self-)reflaction or validation of what was extracted?

It is also not clear (i) the kind of documents used in this process (paper, patents, books, ....), (ii) if different agents/LLMs were used to extract information from these different documents, (iii) the relation between the different documents and the information extracted (such as (canonical name, HS code, ...)

The architecture presented in Fig 1 is only a list of buzzwords. It is not clear how the piaces are connects and what are the contributions of having used agents in the process.
\color{black}

\textbf{Response}
Section 3 explains our agentic approach, which uses LLM agents to construct the knowledge hypergraphs.
Details of this approach are presented in Section 3.4, and we feel this is a fit for AAMAS due to its agent architecture.
The datasets used by the LLMs are clarified in Section 3.5.
We will elaborate on this in the final paper.

